\section{Research Questions}\label{sec:rqs}

The two events share a warning period but end differently: one with an enforced cutoff that breaks CI, the other with a cutoff that was announced and then withdrawn.
We use this contrast to ask how projects respond.

\textbf{RQ1: When do projects migrate away from a deprecated component?}
We place each migration on the timeline of its deprecation: before anything broke, during the brownouts, after the cutoff, or never.
This tells whether warnings alone move projects, or only breakage does.

\textbf{RQ2: How does the way a deprecation ends affect the speed and extent of migration?}
We compare the time to migration under the enforced and the withdrawn cutoff, overall and within repositories exposed to both.

\textbf{RQ3: Which project characteristics are associated with timely migration?}
We relate migrating before any breakage to the size, activity and dependency tooling (Dependabot, Renovate) of a project.

\textbf{RQ4: Who performs the migrations?}
We classify the author of each migration as a dependency bot, a core maintainer or an external contributor.

\textbf{RQ5: How do projects migrate, and what problems do they meet?}
We code each migration by the kind of change, what triggered it, and the problems visible in the commits.
